\section{Status of the analytical program and falsifiable tests}
\label{sec:status-tests}

The construction yields a precise number from a precise branch quotient.  The electroweak comparison is meaningful only after choosing the pole-mass observable.  The project advances when the following calculations are completed.

\begin{enumerate}
\item Derive the ordered assignment \((J_W,J_Z)=(3/4,2)\) from the gauge-Higgs sector or from a boundary/compactification problem.
\item Derive the two-channel determinant Eq.~\eqref{eq:det-completion} from a dynamical model.
\item Determine whether the negative branch maps to the Higgs instability or to another auxiliary sector.
\item Update the pole-mass comparison with current source-audited inputs and uncertainty propagation.
\item Specify the global form of the gauge group in any proposed topological completion.
\end{enumerate}

The falsifiable tests are correspondingly direct.  A source-audited pole-mass ratio that decisively departs from Eq.~\eqref{eq:sdvdef} would weaken the central observation.  A proof that no gauge-Higgs or boundary construction can produce Eq.~\eqref{eq:Jassignment} under the required assumptions would block the electroweak interpretation.  A dynamical derivation that produces a different ordered pair or a different branch quotient would redirect the project toward a broader class of secular mass relations.

\paragraph*{Section status.}  This section states the program's current logical dependencies.
